% Options for packages loaded elsewhere
% Options for packages loaded elsewhere
\PassOptionsToPackage{unicode}{hyperref}
\PassOptionsToPackage{hyphens}{url}
\PassOptionsToPackage{dvipsnames,svgnames,x11names}{xcolor}
%
\documentclass[
  portuguese,
  letterpaper,
  DIV=11,
  numbers=noendperiod]{scrartcl}
\usepackage{xcolor}
\usepackage{amsmath,amssymb}
\setcounter{secnumdepth}{5}
\usepackage{iftex}
\ifPDFTeX
  \usepackage[T1]{fontenc}
  \usepackage[utf8]{inputenc}
  \usepackage{textcomp} % provide euro and other symbols
\else % if luatex or xetex
  \usepackage{unicode-math} % this also loads fontspec
  \defaultfontfeatures{Scale=MatchLowercase}
  \defaultfontfeatures[\rmfamily]{Ligatures=TeX,Scale=1}
\fi
\usepackage{lmodern}
\ifPDFTeX\else
  % xetex/luatex font selection
\fi
% Use upquote if available, for straight quotes in verbatim environments
\IfFileExists{upquote.sty}{\usepackage{upquote}}{}
\IfFileExists{microtype.sty}{% use microtype if available
  \usepackage[]{microtype}
  \UseMicrotypeSet[protrusion]{basicmath} % disable protrusion for tt fonts
}{}
\makeatletter
\@ifundefined{KOMAClassName}{% if non-KOMA class
  \IfFileExists{parskip.sty}{%
    \usepackage{parskip}
  }{% else
    \setlength{\parindent}{0pt}
    \setlength{\parskip}{6pt plus 2pt minus 1pt}}
}{% if KOMA class
  \KOMAoptions{parskip=half}}
\makeatother
% Make \paragraph and \subparagraph free-standing
\makeatletter
\ifx\paragraph\undefined\else
  \let\oldparagraph\paragraph
  \renewcommand{\paragraph}{
    \@ifstar
      \xxxParagraphStar
      \xxxParagraphNoStar
  }
  \newcommand{\xxxParagraphStar}[1]{\oldparagraph*{#1}\mbox{}}
  \newcommand{\xxxParagraphNoStar}[1]{\oldparagraph{#1}\mbox{}}
\fi
\ifx\subparagraph\undefined\else
  \let\oldsubparagraph\subparagraph
  \renewcommand{\subparagraph}{
    \@ifstar
      \xxxSubParagraphStar
      \xxxSubParagraphNoStar
  }
  \newcommand{\xxxSubParagraphStar}[1]{\oldsubparagraph*{#1}\mbox{}}
  \newcommand{\xxxSubParagraphNoStar}[1]{\oldsubparagraph{#1}\mbox{}}
\fi
\makeatother

\usepackage{color}
\usepackage{fancyvrb}
\newcommand{\VerbBar}{|}
\newcommand{\VERB}{\Verb[commandchars=\\\{\}]}
\DefineVerbatimEnvironment{Highlighting}{Verbatim}{commandchars=\\\{\}}
% Add ',fontsize=\small' for more characters per line
\usepackage{framed}
\definecolor{shadecolor}{RGB}{241,243,245}
\newenvironment{Shaded}{\begin{snugshade}}{\end{snugshade}}
\newcommand{\AlertTok}[1]{\textcolor[rgb]{0.68,0.00,0.00}{#1}}
\newcommand{\AnnotationTok}[1]{\textcolor[rgb]{0.37,0.37,0.37}{#1}}
\newcommand{\AttributeTok}[1]{\textcolor[rgb]{0.40,0.45,0.13}{#1}}
\newcommand{\BaseNTok}[1]{\textcolor[rgb]{0.68,0.00,0.00}{#1}}
\newcommand{\BuiltInTok}[1]{\textcolor[rgb]{0.00,0.23,0.31}{#1}}
\newcommand{\CharTok}[1]{\textcolor[rgb]{0.13,0.47,0.30}{#1}}
\newcommand{\CommentTok}[1]{\textcolor[rgb]{0.37,0.37,0.37}{#1}}
\newcommand{\CommentVarTok}[1]{\textcolor[rgb]{0.37,0.37,0.37}{\textit{#1}}}
\newcommand{\ConstantTok}[1]{\textcolor[rgb]{0.56,0.35,0.01}{#1}}
\newcommand{\ControlFlowTok}[1]{\textcolor[rgb]{0.00,0.23,0.31}{\textbf{#1}}}
\newcommand{\DataTypeTok}[1]{\textcolor[rgb]{0.68,0.00,0.00}{#1}}
\newcommand{\DecValTok}[1]{\textcolor[rgb]{0.68,0.00,0.00}{#1}}
\newcommand{\DocumentationTok}[1]{\textcolor[rgb]{0.37,0.37,0.37}{\textit{#1}}}
\newcommand{\ErrorTok}[1]{\textcolor[rgb]{0.68,0.00,0.00}{#1}}
\newcommand{\ExtensionTok}[1]{\textcolor[rgb]{0.00,0.23,0.31}{#1}}
\newcommand{\FloatTok}[1]{\textcolor[rgb]{0.68,0.00,0.00}{#1}}
\newcommand{\FunctionTok}[1]{\textcolor[rgb]{0.28,0.35,0.67}{#1}}
\newcommand{\ImportTok}[1]{\textcolor[rgb]{0.00,0.46,0.62}{#1}}
\newcommand{\InformationTok}[1]{\textcolor[rgb]{0.37,0.37,0.37}{#1}}
\newcommand{\KeywordTok}[1]{\textcolor[rgb]{0.00,0.23,0.31}{\textbf{#1}}}
\newcommand{\NormalTok}[1]{\textcolor[rgb]{0.00,0.23,0.31}{#1}}
\newcommand{\OperatorTok}[1]{\textcolor[rgb]{0.37,0.37,0.37}{#1}}
\newcommand{\OtherTok}[1]{\textcolor[rgb]{0.00,0.23,0.31}{#1}}
\newcommand{\PreprocessorTok}[1]{\textcolor[rgb]{0.68,0.00,0.00}{#1}}
\newcommand{\RegionMarkerTok}[1]{\textcolor[rgb]{0.00,0.23,0.31}{#1}}
\newcommand{\SpecialCharTok}[1]{\textcolor[rgb]{0.37,0.37,0.37}{#1}}
\newcommand{\SpecialStringTok}[1]{\textcolor[rgb]{0.13,0.47,0.30}{#1}}
\newcommand{\StringTok}[1]{\textcolor[rgb]{0.13,0.47,0.30}{#1}}
\newcommand{\VariableTok}[1]{\textcolor[rgb]{0.07,0.07,0.07}{#1}}
\newcommand{\VerbatimStringTok}[1]{\textcolor[rgb]{0.13,0.47,0.30}{#1}}
\newcommand{\WarningTok}[1]{\textcolor[rgb]{0.37,0.37,0.37}{\textit{#1}}}

\usepackage{longtable,booktabs,array}
\usepackage{calc} % for calculating minipage widths
% Correct order of tables after \paragraph or \subparagraph
\usepackage{etoolbox}
\makeatletter
\patchcmd\longtable{\par}{\if@noskipsec\mbox{}\fi\par}{}{}
\makeatother
% Allow footnotes in longtable head/foot
\IfFileExists{footnotehyper.sty}{\usepackage{footnotehyper}}{\usepackage{footnote}}
\makesavenoteenv{longtable}
\usepackage{graphicx}
\makeatletter
\newsavebox\pandoc@box
\newcommand*\pandocbounded[1]{% scales image to fit in text height/width
  \sbox\pandoc@box{#1}%
  \Gscale@div\@tempa{\textheight}{\dimexpr\ht\pandoc@box+\dp\pandoc@box\relax}%
  \Gscale@div\@tempb{\linewidth}{\wd\pandoc@box}%
  \ifdim\@tempb\p@<\@tempa\p@\let\@tempa\@tempb\fi% select the smaller of both
  \ifdim\@tempa\p@<\p@\scalebox{\@tempa}{\usebox\pandoc@box}%
  \else\usebox{\pandoc@box}%
  \fi%
}
% Set default figure placement to htbp
\def\fps@figure{htbp}
\makeatother



\ifLuaTeX
\usepackage[bidi=basic]{babel}
\else
\usepackage[bidi=default]{babel}
\fi
% get rid of language-specific shorthands (see #6817):
\let\LanguageShortHands\languageshorthands
\def\languageshorthands#1{}


\setlength{\emergencystretch}{3em} % prevent overfull lines

\providecommand{\tightlist}{%
  \setlength{\itemsep}{0pt}\setlength{\parskip}{0pt}}



 


\usepackage{booktabs}
\usepackage{longtable}
\usepackage{array}
\usepackage{multirow}
\usepackage{wrapfig}
\usepackage{float}
\usepackage{colortbl}
\usepackage{pdflscape}
\usepackage{tabu}
\usepackage{threeparttable}
\usepackage{threeparttablex}
\usepackage[normalem]{ulem}
\usepackage{makecell}
\usepackage{xcolor}
\usepackage{siunitx}

    \newcolumntype{d}{S[
      table-align-text-before=false,
      table-align-text-after=false,
      input-symbols={-,\*+()}
    ]}
  
\KOMAoption{captions}{tableheading}
\makeatletter
\@ifpackageloaded{caption}{}{\usepackage{caption}}
\AtBeginDocument{%
\ifdefined\contentsname
  \renewcommand*\contentsname{Índice}
\else
  \newcommand\contentsname{Índice}
\fi
\ifdefined\listfigurename
  \renewcommand*\listfigurename{Lista de Figuras}
\else
  \newcommand\listfigurename{Lista de Figuras}
\fi
\ifdefined\listtablename
  \renewcommand*\listtablename{Lista de Tabelas}
\else
  \newcommand\listtablename{Lista de Tabelas}
\fi
\ifdefined\figurename
  \renewcommand*\figurename{Figura}
\else
  \newcommand\figurename{Figura}
\fi
\ifdefined\tablename
  \renewcommand*\tablename{Tabela}
\else
  \newcommand\tablename{Tabela}
\fi
}
\@ifpackageloaded{float}{}{\usepackage{float}}
\floatstyle{ruled}
\@ifundefined{c@chapter}{\newfloat{codelisting}{h}{lop}}{\newfloat{codelisting}{h}{lop}[chapter]}
\floatname{codelisting}{Listagem}
\newcommand*\listoflistings{\listof{codelisting}{Lista de Listagens}}
\makeatother
\makeatletter
\makeatother
\makeatletter
\@ifpackageloaded{caption}{}{\usepackage{caption}}
\@ifpackageloaded{subcaption}{}{\usepackage{subcaption}}
\makeatother
\usepackage{bookmark}
\IfFileExists{xurl.sty}{\usepackage{xurl}}{} % add URL line breaks if available
\urlstyle{same}
\hypersetup{
  pdftitle={E10.2 --- Renda e Democracia},
  pdflang={pt},
  colorlinks=true,
  linkcolor={blue},
  filecolor={Maroon},
  citecolor={Blue},
  urlcolor={Blue},
  pdfcreator={LaTeX via pandoc}}


\title{E10.2 --- Renda e Democracia}
\usepackage{etoolbox}
\makeatletter
\providecommand{\subtitle}[1]{% add subtitle to \maketitle
  \apptocmd{\@title}{\par {\large #1 \par}}{}{}
}
\makeatother
\subtitle{Solução --- Regressão com Dados em Painel}
\author{}
\date{2026-03-28}
\begin{document}
\maketitle

\renewcommand*\contentsname{Índice}
{
\hypersetup{linkcolor=}
\setcounter{tocdepth}{3}
\tableofcontents
}

\section{Introdução}\label{introduuxe7uxe3o}

Este documento apresenta a solução para o exercício E10.2 de Stock \&
Watson (2020), que investiga a relação entre renda per capita e
democracia utilizando um painel de 195 países para os anos de 1960 a
2000 (em intervalos de 5 anos).

A pergunta central é: \textbf{a democracia é um bem normal?} Ou seja, os
cidadãos demandam mais liberdade política à medida que sua renda cresce?

\begin{center}\rule{0.5\linewidth}{0.5pt}\end{center}

\section{Pacotes e Dados}\label{pacotes-e-dados}

\begin{Shaded}
\begin{Highlighting}[]
\FunctionTok{library}\NormalTok{(tidyverse)}
\FunctionTok{library}\NormalTok{(skimr)}
\FunctionTok{library}\NormalTok{(fixest)}
\FunctionTok{library}\NormalTok{(modelsummary)}
\FunctionTok{library}\NormalTok{(haven)      }\CommentTok{\# para leitura de arquivos .dta (Stata)}
\FunctionTok{library}\NormalTok{(kableExtra)}
\FunctionTok{library}\NormalTok{(readxl)}

\CommentTok{\# Carregue o arquivo de dados:}
\CommentTok{\# df \textless{}{-} read\_dta("Income\_Democracy.dta")}
\CommentTok{\#}
\CommentTok{\# Para fins de demonstração, a estrutura esperada dos dados é:}
\CommentTok{\#   {-} dem\_ind     : índice de democracia (0 a 1)}
\CommentTok{\#   {-} log\_gdppc   : log da renda per capita (defasado 5 anos)}
\CommentTok{\#   {-} country     : identificador do país}
\CommentTok{\#   {-} year        : ano da observação}
\CommentTok{\#   {-} dem\_ind, log\_gdppc, age\_1, age\_2, age\_3, age\_4, educ, age\_median (controles)}
\CommentTok{\#}
\CommentTok{\# Substitua o caminho abaixo pelo caminho local do arquivo:}
\NormalTok{df }\OtherTok{\textless{}{-}} \FunctionTok{read\_xlsx}\NormalTok{(here}\SpecialCharTok{::}\FunctionTok{here}\NormalTok{(}\StringTok{"labs"}\NormalTok{,}\StringTok{"income\_democracy"}\NormalTok{,}\StringTok{"income\_democracy.xlsx"}\NormalTok{))}
\end{Highlighting}
\end{Shaded}

\begin{quote}
\textbf{Nota:} Os blocos de código abaixo assumem que o objeto
\texttt{df} foi carregado com o arquivo \texttt{Income\_Democracy.dta}
disponível no site do livro. Substitua o caminho conforme necessário.
\end{quote}

\begin{center}\rule{0.5\linewidth}{0.5pt}\end{center}

\section{(a) Painel Balanceado?}\label{a-painel-balanceado}

\begin{Shaded}
\begin{Highlighting}[]
\CommentTok{\# Verificar se o painel é balanceado}
\NormalTok{df }\SpecialCharTok{|\textgreater{}}
  \FunctionTok{count}\NormalTok{(country) }\SpecialCharTok{|\textgreater{}}
  \FunctionTok{summarise}\NormalTok{(}
    \AttributeTok{n\_paises        =} \FunctionTok{n}\NormalTok{(),}
    \AttributeTok{obs\_min         =} \FunctionTok{min}\NormalTok{(n),}
    \AttributeTok{obs\_max         =} \FunctionTok{max}\NormalTok{(n),}
    \AttributeTok{painel\_balanceado =} \FunctionTok{if\_else}\NormalTok{(obs\_min }\SpecialCharTok{==}\NormalTok{ obs\_max, }\StringTok{"Sim"}\NormalTok{, }\StringTok{"Não"}\NormalTok{)}
\NormalTok{  ) }\SpecialCharTok{|\textgreater{}}
  \FunctionTok{kbl}\NormalTok{(}\AttributeTok{caption =} \StringTok{"Verificação de Painel Balanceado"}\NormalTok{,}
      \AttributeTok{col.names =} \FunctionTok{c}\NormalTok{(}\StringTok{"Nº de Países"}\NormalTok{, }\StringTok{"Mín. Obs."}\NormalTok{, }\StringTok{"Máx. Obs."}\NormalTok{, }\StringTok{"Balanceado?"}\NormalTok{)) }\SpecialCharTok{|\textgreater{}}
  \FunctionTok{kable\_styling}\NormalTok{(}\AttributeTok{bootstrap\_options =} \FunctionTok{c}\NormalTok{(}\StringTok{"striped"}\NormalTok{, }\StringTok{"hover"}\NormalTok{), }\AttributeTok{full\_width =} \ConstantTok{FALSE}\NormalTok{)}
\end{Highlighting}
\end{Shaded}

\begin{table}
\centering
\caption{\label{tab:painel-balanceado}Verificação de Painel Balanceado}
\centering
\begin{tabular}[t]{r|r|r|l}
\hline
Nº de Países & Mín. Obs. & Máx. Obs. & Balanceado?\\
\hline
195 & 1 & 9 & Não\\
\hline
\end{tabular}
\end{table}

Um painel é \textbf{balanceado} quando todos os países possuem o mesmo
número de observações ao longo do tempo. Caso contrário, é \textbf{não
balanceado}, o que pode ocorrer por dados faltantes em alguns
países/anos.

\begin{center}\rule{0.5\linewidth}{0.5pt}\end{center}

\section{\texorpdfstring{(b) Estatísticas Descritivas de
\texttt{Dem\_ind}}{(b) Estatísticas Descritivas de Dem\_ind}}\label{b-estatuxedsticas-descritivas-de-dem_ind}

\subsection{(b.i--b.iv) Tabela de Estatísticas
Descritivas}\label{b.ib.iv-tabela-de-estatuxedsticas-descritivas}

\begin{Shaded}
\begin{Highlighting}[]
\CommentTok{\# Estatísticas descritivas com skimr}
\NormalTok{skim\_resultado }\OtherTok{\textless{}{-}}\NormalTok{ df }\SpecialCharTok{|\textgreater{}}
  \FunctionTok{select}\NormalTok{(dem\_ind, log\_gdppc) }\SpecialCharTok{|\textgreater{}}
  \FunctionTok{skim}\NormalTok{()}

\CommentTok{\# Tabela formatada com percentis customizados}
\NormalTok{desc\_table }\OtherTok{\textless{}{-}}\NormalTok{ df }\SpecialCharTok{|\textgreater{}}
  \FunctionTok{summarise}\NormalTok{(}
\NormalTok{    Variável     }\OtherTok{=} \StringTok{"dem\_ind"}\NormalTok{,}
    \AttributeTok{N            =} \FunctionTok{sum}\NormalTok{(}\SpecialCharTok{!}\FunctionTok{is.na}\NormalTok{(dem\_ind)),}
\NormalTok{    Mínimo       }\OtherTok{=} \FunctionTok{min}\NormalTok{(dem\_ind, }\AttributeTok{na.rm =} \ConstantTok{TRUE}\NormalTok{),}
\NormalTok{    Máximo       }\OtherTok{=} \FunctionTok{max}\NormalTok{(dem\_ind, }\AttributeTok{na.rm =} \ConstantTok{TRUE}\NormalTok{),}
\NormalTok{    Média        }\OtherTok{=} \FunctionTok{mean}\NormalTok{(dem\_ind, }\AttributeTok{na.rm =} \ConstantTok{TRUE}\NormalTok{),}
    \StringTok{\textasciigrave{}}\AttributeTok{Desvio Padrão}\StringTok{\textasciigrave{}} \OtherTok{=} \FunctionTok{sd}\NormalTok{(dem\_ind, }\AttributeTok{na.rm =} \ConstantTok{TRUE}\NormalTok{),}
    \StringTok{\textasciigrave{}}\AttributeTok{P10}\StringTok{\textasciigrave{}}        \OtherTok{=} \FunctionTok{quantile}\NormalTok{(dem\_ind, }\FloatTok{0.10}\NormalTok{, }\AttributeTok{na.rm =} \ConstantTok{TRUE}\NormalTok{),}
    \StringTok{\textasciigrave{}}\AttributeTok{P25}\StringTok{\textasciigrave{}}        \OtherTok{=} \FunctionTok{quantile}\NormalTok{(dem\_ind, }\FloatTok{0.25}\NormalTok{, }\AttributeTok{na.rm =} \ConstantTok{TRUE}\NormalTok{),}
    \StringTok{\textasciigrave{}}\AttributeTok{P50}\StringTok{\textasciigrave{}}        \OtherTok{=} \FunctionTok{quantile}\NormalTok{(dem\_ind, }\FloatTok{0.50}\NormalTok{, }\AttributeTok{na.rm =} \ConstantTok{TRUE}\NormalTok{),}
    \StringTok{\textasciigrave{}}\AttributeTok{P75}\StringTok{\textasciigrave{}}        \OtherTok{=} \FunctionTok{quantile}\NormalTok{(dem\_ind, }\FloatTok{0.75}\NormalTok{, }\AttributeTok{na.rm =} \ConstantTok{TRUE}\NormalTok{),}
    \StringTok{\textasciigrave{}}\AttributeTok{P90}\StringTok{\textasciigrave{}}        \OtherTok{=} \FunctionTok{quantile}\NormalTok{(dem\_ind, }\FloatTok{0.90}\NormalTok{, }\AttributeTok{na.rm =} \ConstantTok{TRUE}\NormalTok{)}
\NormalTok{  )}

\NormalTok{desc\_table }\SpecialCharTok{|\textgreater{}}
  \FunctionTok{kbl}\NormalTok{(}\AttributeTok{digits =} \DecValTok{3}\NormalTok{,}
      \AttributeTok{caption =} \StringTok{"Tabela 1 — Estatísticas Descritivas do Índice de Democracia (Dem\_ind)"}\NormalTok{) }\SpecialCharTok{|\textgreater{}}
  \FunctionTok{kable\_styling}\NormalTok{(}\AttributeTok{bootstrap\_options =} \FunctionTok{c}\NormalTok{(}\StringTok{"striped"}\NormalTok{, }\StringTok{"hover"}\NormalTok{), }\AttributeTok{full\_width =} \ConstantTok{FALSE}\NormalTok{)}
\end{Highlighting}
\end{Shaded}

\begin{table}
\centering
\caption{\label{tab:estatisticas-descritivas}Tabela 1 — Estatísticas Descritivas do Índice de Democracia (Dem_ind)}
\centering
\begin{tabular}[t]{l|r|r|r|r|r|r|r|r|r|r}
\hline
Variável & N & Mínimo & Máximo & Média & Desvio Padrão & P10 & P25 & P50 & P75 & P90\\
\hline
dem\_ind & 1266 & 0 & 1 & 0.499 & 0.371 & 0 & 0.167 & 0.5 & 0.833 & 1\\
\hline
\end{tabular}
\end{table}

\begin{Shaded}
\begin{Highlighting}[]
\CommentTok{\# Valores para EUA e Líbia}
\NormalTok{paises\_interesse }\OtherTok{\textless{}{-}}\NormalTok{ df }\SpecialCharTok{|\textgreater{}}
  \FunctionTok{filter}\NormalTok{(country }\SpecialCharTok{\%in\%} \FunctionTok{c}\NormalTok{(}\StringTok{"United States"}\NormalTok{, }\StringTok{"Libya"}\NormalTok{)) }\SpecialCharTok{|\textgreater{}}
  \FunctionTok{group\_by}\NormalTok{(country) }\SpecialCharTok{|\textgreater{}}
  \FunctionTok{summarise}\NormalTok{(}
    \StringTok{\textasciigrave{}}\AttributeTok{Dem\_ind em 2000}\StringTok{\textasciigrave{}} \OtherTok{=}\NormalTok{ dem\_ind[year }\SpecialCharTok{==} \DecValTok{2000}\NormalTok{],}
    \StringTok{\textasciigrave{}}\AttributeTok{Média (todos os anos)}\StringTok{\textasciigrave{}} \OtherTok{=} \FunctionTok{mean}\NormalTok{(dem\_ind, }\AttributeTok{na.rm =} \ConstantTok{TRUE}\NormalTok{)}
\NormalTok{  )}

\NormalTok{paises\_interesse }\SpecialCharTok{|\textgreater{}}
  \FunctionTok{kbl}\NormalTok{(}\AttributeTok{digits =} \DecValTok{3}\NormalTok{,}
      \AttributeTok{caption =} \StringTok{"Dem\_ind para EUA e Líbia"}\NormalTok{) }\SpecialCharTok{|\textgreater{}}
  \FunctionTok{kable\_styling}\NormalTok{(}\AttributeTok{bootstrap\_options =} \FunctionTok{c}\NormalTok{(}\StringTok{"striped"}\NormalTok{, }\StringTok{"hover"}\NormalTok{), }\AttributeTok{full\_width =} \ConstantTok{FALSE}\NormalTok{)}
\end{Highlighting}
\end{Shaded}

\begin{table}
\centering
\caption{\label{tab:paises-especificos}Dem_ind para EUA e Líbia}
\centering
\begin{tabular}[t]{l|r|r}
\hline
country & Dem\_ind em 2000 & Média (todos os anos)\\
\hline
Libya & 0 & 0.109\\
\hline
United States & 1 & 0.986\\
\hline
\end{tabular}
\end{table}

\begin{Shaded}
\begin{Highlighting}[]
\CommentTok{\# Países por faixa de Dem\_ind médio}
\NormalTok{df }\SpecialCharTok{|\textgreater{}}
  \FunctionTok{group\_by}\NormalTok{(country) }\SpecialCharTok{|\textgreater{}}
  \FunctionTok{summarise}\NormalTok{(}\AttributeTok{media\_dem =} \FunctionTok{mean}\NormalTok{(dem\_ind, }\AttributeTok{na.rm =} \ConstantTok{TRUE}\NormalTok{)) }\SpecialCharTok{|\textgreater{}}
  \FunctionTok{filter}\NormalTok{(media\_dem }\SpecialCharTok{\textgreater{}} \FloatTok{0.95} \SpecialCharTok{|}\NormalTok{ media\_dem }\SpecialCharTok{\textless{}} \FloatTok{0.10} \SpecialCharTok{|} \FunctionTok{between}\NormalTok{(media\_dem, }\FloatTok{0.3}\NormalTok{, }\FloatTok{0.7}\NormalTok{)) }\SpecialCharTok{|\textgreater{}}
  \FunctionTok{mutate}\NormalTok{(}\AttributeTok{faixa =} \FunctionTok{case\_when}\NormalTok{(}
\NormalTok{    media\_dem }\SpecialCharTok{\textgreater{}} \FloatTok{0.95}            \SpecialCharTok{\textasciitilde{}} \StringTok{"\textgreater{} 0,95"}\NormalTok{,}
\NormalTok{    media\_dem }\SpecialCharTok{\textless{}} \FloatTok{0.10}            \SpecialCharTok{\textasciitilde{}} \StringTok{"\textless{} 0,10"}\NormalTok{,}
    \FunctionTok{between}\NormalTok{(media\_dem, }\FloatTok{0.3}\NormalTok{, }\FloatTok{0.7}\NormalTok{) }\SpecialCharTok{\textasciitilde{}} \StringTok{"Entre 0,3 e 0,7"}
\NormalTok{  )) }\SpecialCharTok{|\textgreater{}}
  \FunctionTok{group\_by}\NormalTok{(faixa) }\SpecialCharTok{|\textgreater{}}
  \FunctionTok{slice\_head}\NormalTok{(}\AttributeTok{n =} \DecValTok{5}\NormalTok{) }\SpecialCharTok{|\textgreater{}}
  \FunctionTok{arrange}\NormalTok{(faixa, }\FunctionTok{desc}\NormalTok{(media\_dem)) }\SpecialCharTok{|\textgreater{}}
  \FunctionTok{kbl}\NormalTok{(}\AttributeTok{digits =} \DecValTok{3}\NormalTok{,}
      \AttributeTok{col.names =} \FunctionTok{c}\NormalTok{(}\StringTok{"País"}\NormalTok{, }\StringTok{"Média de Dem\_ind"}\NormalTok{, }\StringTok{"Faixa"}\NormalTok{),}
      \AttributeTok{caption =} \StringTok{"Países selecionados por faixa de Dem\_ind médio"}\NormalTok{) }\SpecialCharTok{|\textgreater{}}
  \FunctionTok{kable\_styling}\NormalTok{(}\AttributeTok{bootstrap\_options =} \FunctionTok{c}\NormalTok{(}\StringTok{"striped"}\NormalTok{, }\StringTok{"hover"}\NormalTok{), }\AttributeTok{full\_width =} \ConstantTok{FALSE}\NormalTok{)}
\end{Highlighting}
\end{Shaded}

\begin{table}
\centering
\caption{\label{tab:paises-especificos}Países selecionados por faixa de Dem_ind médio}
\centering
\begin{tabular}[t]{l|r|l}
\hline
País & Média de Dem\_ind & Faixa\\
\hline
Afghanistan & 0.097 & < 0,10\\
\hline
China & 0.080 & < 0,10\\
\hline
Angola & 0.067 & < 0,10\\
\hline
Brunei & 0.056 & < 0,10\\
\hline
Burundi & 0.048 & < 0,10\\
\hline
Australia & 1.000 & > 0,95\\
\hline
Barbados & 1.000 & > 0,95\\
\hline
Belgium & 1.000 & > 0,95\\
\hline
Belize & 1.000 & > 0,95\\
\hline
Austria & 0.993 & > 0,95\\
\hline
Argentina & 0.666 & Entre 0,3 e 0,7\\
\hline
Antigua & 0.556 & Entre 0,3 e 0,7\\
\hline
Bolivia & 0.551 & Entre 0,3 e 0,7\\
\hline
Bangladesh & 0.533 & Entre 0,3 e 0,7\\
\hline
Armenia & 0.500 & Entre 0,3 e 0,7\\
\hline
\end{tabular}
\end{table}

\begin{center}\rule{0.5\linewidth}{0.5pt}\end{center}

\section{(c--d) Estimações em
Painel}\label{cd-estimauxe7uxf5es-em-painel}

As regressões abaixo seguem a sequência do exercício, adicionando
progressivamente efeitos fixos e controles demográficos.

\begin{Shaded}
\begin{Highlighting}[]
\CommentTok{\# (c) OLS simples com erros clusterizados por país}
\NormalTok{m1 }\OtherTok{\textless{}{-}} \FunctionTok{feols}\NormalTok{(dem\_ind }\SpecialCharTok{\textasciitilde{}}\NormalTok{ log\_gdppc,}
            \AttributeTok{data    =}\NormalTok{ df,}
            \AttributeTok{cluster =} \SpecialCharTok{\textasciitilde{}}\NormalTok{country)}

\CommentTok{\# (d.ii) Efeitos fixos de país}
\NormalTok{m2 }\OtherTok{\textless{}{-}} \FunctionTok{feols}\NormalTok{(dem\_ind }\SpecialCharTok{\textasciitilde{}}\NormalTok{ log\_gdppc }\SpecialCharTok{|}\NormalTok{ country,}
            \AttributeTok{data    =}\NormalTok{ df,}
            \AttributeTok{cluster =} \SpecialCharTok{\textasciitilde{}}\NormalTok{country)}

\CommentTok{\# (d.ii) sem Azerbaijão}
\NormalTok{m3 }\OtherTok{\textless{}{-}} \FunctionTok{feols}\NormalTok{(dem\_ind }\SpecialCharTok{\textasciitilde{}}\NormalTok{ log\_gdppc }\SpecialCharTok{|}\NormalTok{ country,}
            \AttributeTok{data    =}\NormalTok{ df }\SpecialCharTok{|\textgreater{}} \FunctionTok{filter}\NormalTok{(country }\SpecialCharTok{!=} \StringTok{"Azerbaijan"}\NormalTok{),}
            \AttributeTok{cluster =} \SpecialCharTok{\textasciitilde{}}\NormalTok{country)}

\CommentTok{\# (d.v) Efeitos fixos de país e tempo}
\NormalTok{m4 }\OtherTok{\textless{}{-}} \FunctionTok{feols}\NormalTok{(dem\_ind }\SpecialCharTok{\textasciitilde{}}\NormalTok{ log\_gdppc }\SpecialCharTok{|}\NormalTok{ country }\SpecialCharTok{+}\NormalTok{ year,}
            \AttributeTok{data    =}\NormalTok{ df,}
            \AttributeTok{cluster =} \SpecialCharTok{\textasciitilde{}}\NormalTok{country)}

\CommentTok{\# (d.vi) Com controles demográficos + efeitos fixos de país e tempo}
\CommentTok{\# (ajuste os nomes das variáveis conforme o dataset)}
\NormalTok{m5 }\OtherTok{\textless{}{-}} \FunctionTok{feols}\NormalTok{(dem\_ind }\SpecialCharTok{\textasciitilde{}}\NormalTok{ log\_gdppc }\SpecialCharTok{+}\NormalTok{ age\_1 }\SpecialCharTok{+}\NormalTok{ age\_2 }\SpecialCharTok{+}\NormalTok{ age\_3 }\SpecialCharTok{+}\NormalTok{ age\_4 }\SpecialCharTok{+}\NormalTok{ educ }\SpecialCharTok{+}\NormalTok{ age\_median }\SpecialCharTok{|}
\NormalTok{              country }\SpecialCharTok{+}\NormalTok{ year,}
            \AttributeTok{data    =}\NormalTok{ df,}
            \AttributeTok{cluster =} \SpecialCharTok{\textasciitilde{}}\NormalTok{country)}
\end{Highlighting}
\end{Shaded}

\subsection{Tabela de Resultados das
Estimações}\label{tabela-de-resultados-das-estimauxe7uxf5es}

\begin{Shaded}
\begin{Highlighting}[]
\FunctionTok{modelsummary}\NormalTok{(}
  \FunctionTok{list}\NormalTok{(}
    \StringTok{"(c) OLS"}                   \OtherTok{=}\NormalTok{ m1,}
    \StringTok{"(d.ii) EF País"}            \OtherTok{=}\NormalTok{ m2,}
    \StringTok{"(d.iii) Sem Azerbaijão"}    \OtherTok{=}\NormalTok{ m3,}
    \StringTok{"(d.v) EF País + Tempo"}     \OtherTok{=}\NormalTok{ m4,}
    \StringTok{"(d.vi) + Controles Demog."} \OtherTok{=}\NormalTok{ m5}
\NormalTok{  ),}
  \AttributeTok{stars      =} \FunctionTok{c}\NormalTok{(}\StringTok{"*"} \OtherTok{=} \FloatTok{0.1}\NormalTok{, }\StringTok{"**"} \OtherTok{=} \FloatTok{0.05}\NormalTok{, }\StringTok{"***"} \OtherTok{=} \FloatTok{0.01}\NormalTok{),}
  \AttributeTok{coef\_map   =} \FunctionTok{c}\NormalTok{(}
    \StringTok{"log\_gdppc"} \OtherTok{=} \StringTok{"Log(PIB per capita)"}\NormalTok{,}
    \StringTok{"age\_1"}     \OtherTok{=} \StringTok{"Faixa etária 1"}\NormalTok{,}
    \StringTok{"age\_2"}     \OtherTok{=} \StringTok{"Faixa etária 2"}\NormalTok{,}
    \StringTok{"age\_3"}     \OtherTok{=} \StringTok{"Faixa etária 3"}\NormalTok{,}
    \StringTok{"age\_4"}     \OtherTok{=} \StringTok{"Faixa etária 4"}\NormalTok{,}
    \StringTok{"educ"}      \OtherTok{=} \StringTok{"Educação"}\NormalTok{,}
    \StringTok{"age\_median"}\OtherTok{=} \StringTok{"Idade mediana"}
\NormalTok{  ),}
  \AttributeTok{gof\_map    =} \FunctionTok{c}\NormalTok{(}\StringTok{"nobs"}\NormalTok{, }\StringTok{"r.squared"}\NormalTok{, }\StringTok{"adj.r.squared"}\NormalTok{),}
  \AttributeTok{title      =} \StringTok{"Tabela 2 — Efeito da Renda per Capita sobre o Índice de Democracia"}\NormalTok{,}
  \AttributeTok{notes      =} \StringTok{"Erros padrão clusterizados por país entre parênteses. * p\textless{}0,1; ** p\textless{}0,05; *** p\textless{}0,01."}\NormalTok{,}
  \AttributeTok{output     =} \StringTok{"kableExtra"}
\NormalTok{) }\SpecialCharTok{|\textgreater{}}
  \FunctionTok{add\_header\_above}\NormalTok{(}\FunctionTok{c}\NormalTok{(}
    \StringTok{" "} \OtherTok{=} \DecValTok{1}\NormalTok{,}
    \StringTok{"Sem Efeitos Fixos"} \OtherTok{=} \DecValTok{1}\NormalTok{,}
    \StringTok{"EF País"} \OtherTok{=} \DecValTok{2}\NormalTok{,}
    \StringTok{"EF País + Tempo"} \OtherTok{=} \DecValTok{2}
\NormalTok{  )) }\SpecialCharTok{|\textgreater{}}
  \FunctionTok{kable\_styling}\NormalTok{(}\AttributeTok{bootstrap\_options =} \FunctionTok{c}\NormalTok{(}\StringTok{"striped"}\NormalTok{, }\StringTok{"hover"}\NormalTok{),}
                \AttributeTok{latex\_options     =} \FunctionTok{c}\NormalTok{(}\StringTok{"hold\_position"}\NormalTok{, }\StringTok{"scale\_down"}\NormalTok{),}
                \AttributeTok{full\_width        =} \ConstantTok{FALSE}\NormalTok{)}
\end{Highlighting}
\end{Shaded}

\begin{table}[!h]
\centering\centering
\caption{\label{tab:tabela-resultados}Tabela 2 — Efeito da Renda per Capita sobre o Índice de Democracia}
\centering
\resizebox{\ifdim\width>\linewidth\linewidth\else\width\fi}{!}{
\begin{tabular}[t]{lccccc}
\toprule
\multicolumn{1}{c}{ } & \multicolumn{1}{c}{Sem Efeitos Fixos} & \multicolumn{2}{c}{EF País} & \multicolumn{2}{c}{EF País + Tempo} \\
\cmidrule(l{3pt}r{3pt}){2-2} \cmidrule(l{3pt}r{3pt}){3-4} \cmidrule(l{3pt}r{3pt}){5-6}
  & (c) OLS & (d.ii) EF País & (d.iii) Sem Azerbaijão & (d.v) EF País + Tempo & (d.vi) + Controles Demog.\\
\midrule
Log(PIB per capita) & \num{0.236}*** & \num{0.084}*** & \num{0.084}*** & \num{0.054} & \num{0.029}\\
 & (\num{0.012}) & (\num{0.032}) & (\num{0.032}) & (\num{0.042}) & (\num{0.054})\\
Faixa etária 1 &  &  &  &  & \num{-1.171}\\
 &  &  &  &  & (\num{1.915})\\
Faixa etária 2 &  &  &  &  & \num{-1.718}\\
 &  &  &  &  & (\num{1.534})\\
Faixa etária 3 &  &  &  &  & \num{-3.508}**\\
 &  &  &  &  & (\num{1.618})\\
Faixa etária 4 &  &  &  &  & \num{-0.895}\\
 &  &  &  &  & (\num{1.641})\\
Educação &  &  &  &  & \num{-0.003}\\
 &  &  &  &  & (\num{0.023})\\
Idade mediana &  &  &  &  & \num{-0.001}\\
 &  &  &  &  & (\num{0.026})\\
\midrule
Num.Obs. & \num{958} & \num{937} & \num{937} & \num{937} & \num{676}\\
R2 & \num{0.438} & \num{0.729} & \num{0.729} & \num{0.756} & \num{0.732}\\
R2 Adj. & \num{0.438} & \num{0.686} & \num{0.686} & \num{0.714} & \num{0.683}\\
\bottomrule
\multicolumn{6}{l}{\rule{0pt}{1em}* p $<$ 0.1, ** p $<$ 0.05, *** p $<$ 0.01}\\
\multicolumn{6}{l}{\rule{0pt}{1em}Erros padrão clusterizados por país entre parênteses. * p<0,1; ** p<0,05; *** p<0,01.}\\
\end{tabular}}
\end{table}

\begin{center}\rule{0.5\linewidth}{0.5pt}\end{center}

\section{(e) Conclusões}\label{e-conclusuxf5es}

Com base nas estimações acima, podemos tirar as seguintes conclusões
sobre os efeitos da renda sobre a democracia:

\begin{enumerate}
\def\labelenumi{\arabic{enumi}.}
\item
  \textbf{Sem efeitos fixos (Modelo 1):} O coeficiente de
  \texttt{log\_gdppc} é positivo e estatisticamente significativo,
  sugerindo uma associação positiva entre renda e democracia.
\item
  \textbf{Com efeitos fixos de país (Modelo 2):} Ao controlar por
  características não observadas e invariantes no tempo de cada país
  (como história, cultura, instituições), o coeficiente se altera ---
  tipicamente diminuindo ---, indicando que parte da correlação anterior
  era espúria.
\item
  \textbf{Excluindo o Azerbaijão (Modelo 3):} Caso o resultado mude de
  forma expressiva, isso sugere que o Azerbaijão é um \emph{outlier}
  influente na estimação.
\item
  \textbf{Com efeitos fixos de país e tempo (Modelo 4):} Controlando
  também por choques globais comuns a todos os países em cada período
  (como crises internacionais ou ondas de democratização), a estimativa
  causal do efeito da renda se torna mais confiável.
\item
  \textbf{Com controles demográficos (Modelo 5):} A inclusão de
  variáveis demográficas permite isolar o efeito da renda de mudanças
  estruturais populacionais que podem afetar simultaneamente renda e
  democracia.
\end{enumerate}

No conjunto, a evidência sugere que \textbf{o efeito causal da renda
sobre a democracia é modesto e menos robusto} do que a correlação bruta
indica, especialmente após a inclusão de efeitos fixos. Isso questiona a
hipótese de que a democracia é simplesmente um ``bem normal'' que cresce
automaticamente com a renda.

\begin{center}\rule{0.5\linewidth}{0.5pt}\end{center}

\section{Referências}\label{referuxeancias}

\begin{itemize}
\tightlist
\item
  Stock, J. H., \& Watson, M. W. (2020). \emph{Introduction to
  Econometrics} (4ª ed.). Pearson.
\item
  Acemoglu, D., Johnson, S., Robinson, J. A., \& Yared, P. (2008).
  Income and Democracy. \emph{American Economic Review}, 98(3),
  808--842.
\end{itemize}




\end{document}
